\chapter{Evaluation Plan}

\section{Functionality Testing}

To test data ingestion, we will assess the correctness in the parsing of CVs, LinkedIn profiles, and GitHub repositories. System output will be compared against a ground-truth dataset.
For metric inference and scoring, the system will be tested on its ability to infer metrics from a sample job description and compute candidate scores. Test scenarios will be designed to emulate clearly good and bad candidates, both in skill level and role fit.
To assess customisability, core configuration features will be manually tested to confirm expected behaviour.
Finally, to evaluate explainable AI reports, we will manually evaluate them for clarity, correctness, and readability, occasionally reaching out to recruiters for feedback.
\\\\
Success for data ingestion will be defined by a $\geq 80\%$ accuracy in parsing key candidate information (per data source).
For metric inference and scoring, success will be defined by the system correctly identifying at least 80\% of key metrics from job descriptions, and producing scores that align with expected rankings in at least 85\% of test cases.
To measure the success for customisability, all customisation options function correctly.
Success for explainable AI reports will be defined by outputs correctly reflecting underlying score calculations.
It will also be considered successful if recruiters find the reports clear and useful.

\section{Effectiveness Evaluation}

MERIT will also be evaluated for its ability to improve early-stage candidate screening compared to conventional CV-only methods.
A sample dataset of candidate profiles will be prepared, including CVs, LinkedIn profiles, and GitHub repositories.
Then, a subset of candidates will be evaluated by HR experts to establish a “ground truth'' ranking.
\\\\
To provide metrics for rankings, we may use Spearman's rank correlation or Kendall's tau to measure the correlation between MERIT's rankings and expert judgments.
To assess the detection of missions skills or weak areas, we will evaluate whether MERIT appropriately flags data gaps and weak competencies.
For bias mitigation effectiveness, we will compare rankings with and without anonymisation to assess differences based on gender and ethnicity bias.
\\\\
For MERIT to be considered effective, the ranking score should correlate at $\geq 0.7$ with expert assessments.
Atleast $80\%$ of relevant candidates should have their missing or weak skill areas correctly flagged.
Finally, observable disparities in ranking based on demographic factors should be reduced after bias mitigation.

\section{Usability and Qualitative Assessment}

Ease of use will also be essential for MERIT's success.
We will conduct user testing sessions with recruiters and HR professionals to gather qualitative feedback on the system's usability, interface design, and overall user experience.
\\\\
We will measure times taken to upload data and generate results.
We will also track task completion rates for core actions such as adjusting weightings.
Finally, subjective satisfaction will be measured through a questionnaire covering ease of use, clarity of reports, and overall satisfaction.
Success will be achieved with high satisfaction scores (average $\geq 4$ out of 5) on usability questionnaires.
Task completion rates should be $\geq 90\%$ for core actions, and end-to-end screening times should be under 10 minutes for typical candidate sets.

\section{Extending the State of the Art}
We are attempting to extend the state of the art in three key areas: multimodal data integration, customisable scoring, and also explainability and transparency.
These are areas that are currently under-explored in automated technical recruitment systems.
To measure success, MERIT must reliably provide these integrations within reason.
As well as this, recruiters must be able to adjust scorings to match role-specific priorities with different weightings meaningfully affecting results.
Recruiters must also understand the generated reports and understand how each score was produced.
